\paragraph{Area-weighted covering calibration.}
We prove the profile in \eqref{eq:fqe-critical-fidelity-floor} using an
explicit probability measure on the sixteen HS planes. Its only
polynomial positivity input is the following finite exact certificate.

\begin{lemma}[Four-point cubic inequality; exact rational certificate]
\label{lem:fqe-four-point-cubic}
If $X,Y\in\mathbb R^4$ satisfy
$\sum_iX_i=\sum_iY_i=0$ and $\sum_iX_i^2=\sum_iY_i^2=4$, then
\begin{equation}\label{eq:fqe-four-point-cubic}
 \sum_{i,j=1}^4|X_i-Y_j|^3\ge64.
\end{equation}
Equality holds at $X=Y=(-1,-1,1,1)$.
\end{lemma}

\begin{proof}
Put $X=Er$, $Y=Es$, where
\[
 E=\begin{pmatrix}1&1&1\\1&-1&-1\\-1&1&-1\\-1&-1&1\end{pmatrix},
 \qquad E^TE=4I,
\]
so that $\|r\|=\|s\|=1$. Every zero-sum four-vector is represented.
Sort $X$ and $Y$ separately and interleave their entries in increasing
order. There are $\binom84=70$ words with four letters of each kind.
Swapping $X,Y$ and simultaneously negating both vectors reduce these
to exactly twenty-three classes (FP-1). Ties may use either order.
For each word let $z_1,\ldots,z_8$ be its ordered entries and
$g_i=z_{i+1}-z_i\ge0$. The absolute values then give a cubic polynomial
$p=\sum_{ij}|X_i-Y_j|^3/16-4$.
For each of the twenty-three words the supplied rational data give
the identity
\begin{equation}\label{eq:fqe-four-point-cubic-sos}
 p=B_0^TG_0B_0+\sum_{i=1}^7g_i B_i^TG_iB_i
       +(\|r\|^2-1)\ell_r+(\|s\|^2-1)\ell_s.
\end{equation}
The $\ell$ are quadratic polynomials and all eight $G_i$ are positive
definite. FP-2 rebuilds all $210$ coefficients through degree four
in each class and proves positivity of all $184$ Grams using integer
arithmetic. Thus the right side is nonnegative on each ordering cone
and the two spheres, proving the full inequality. The bases and exact
matrix representation are specified in Section~\ref{sec:sup-verifiers}.
The stated equality follows by substitution.
\end{proof}

Let $U=(U_1,U_2)$ be a real orthonormal $4\times2$ frame with plane projector
$\Pi=UU^T$. For each HS plane $P$ use its orthonormal frame
$V_P=(e_k,u)$ and put $c_P=|\det(U^TV_P)|$. The sixteen planes obey
\begin{equation}\label{eq:fqe-area-weighted-fusion}
 \sum_Pc_P^2=\frac83,\qquad
 \sum_Pc_P^2P=\frac89(I+\Pi).
\end{equation}
Indeed, the bivector $\Delta=U_1U_2^T-U_2U_1^T$ satisfies $\Delta\Delta^T=\Delta^T\Delta=\Pi$, and
$\det(U^TV_P)=e_k^T\Delta u$. Fix $k$, put $g=\Delta^Te_k$ (so $g_k=0$,
$|g|^2=\Pi_{kk}$) and sum over the four lines $u$ for this $k$, using $u_j=\pm1/\sqrt3$ ($j\ne k$):
$\sum_u(g\cdot u)^2=\frac43|g|^2$ and
$\sum_u(g\cdot u)^2uu^T=\frac49|g|^2(I-e_ke_k^T)+\frac89(gg^T-\operatorname{diag}(g_j^2))$.
Summing over $k$ with $\sum_kgg^T=\Pi$ and $\sum_k\operatorname{diag}(g_j^2)=\operatorname{diag}(\Pi)$ gives
both identities. FP-3 also checks them without
assuming orthonormality first, using the homogeneous squared area and
its associated matrix. Therefore the explicit weights
$w_P=3c_P^2/8$ satisfy
\begin{equation}\label{eq:fqe-plane-calibration}
 \sum_Pw_P=1,\qquad \sum_Pw_PU^TPU=\frac23I_2,\qquad
 \sum_Pw_Pc_P\ge\frac1{\sqrt3}.
\end{equation}

To prove the last inequality, let $M$ be the magic-basis matrix defined before
\eqref{eq:fqe-null-critical-bloch}. Form the null critical vector
$z=(U_1+iU_2)/\sqrt2$. Its auxiliary product representation
$Mz=a\otimes b$ has unit Bloch vectors $r,s$, as in
\eqref{eq:fqe-null-critical-bloch}.
For the two HS rays $\eta_P=(e_k+iu)/\sqrt2$ and $\bar\eta_P$,
expansion gives
\[
 c_P=\bigl||\langle\eta_P,z\rangle|^2
                 -|\langle\bar\eta_P,z\rangle|^2\bigr|.
\]
In the cube-pair description of the HS rays given after \eqref{eq:fqe-null-critical-bloch}, conjugation reverses both cube vertices.
Indeed $M=\epsilon^{\otimes2}\bar M$, where
$\epsilon=\left(\begin{smallmatrix}0&1\\-1&0\end{smallmatrix}\right)$
(FP-3), and the qubit spin flip reverses the Bloch vector.
Choose the four even-parity cube vertices $e_i$, the rows of $E$.
The sixteen conjugate pairs are represented once by
$(\varepsilon,\eta)=(e_i,-e_j)$, and hence
\[
 c_P=\frac{|e_i\cdot r-e_j\cdot s|}{2\sqrt3}.
\]
The four numbers $X_i=e_i\cdot r$, and likewise $Y_j=e_j\cdot s$,
have sum zero and sum of squares four. Lemma~\ref{lem:fqe-four-point-cubic}
therefore gives
\[
 \sum_Pc_P^3=\frac{\sum_{ij}|X_i-Y_j|^3}{24\sqrt3}
       \ge\frac8{3\sqrt3},
\]
which proves the last calibration inequality.

Now write the real density matrix of any critical state as
$A=U\Lambda U^T$, with $\operatorname{tr}\Lambda=1$ and
$\det\Lambda=(1-\tau)/4$. Rank one is included by completing its
support to any plane. The exact Bures formula gives
\begin{align*}
 \Phi
 &\ge\sum_Pw_P F_{\rm B}(A,P/2)\\
 &=\frac12\operatorname{tr}\!\left(\Lambda\sum_Pw_PU^TPU\right)
                  +\sqrt{\det\Lambda}\sum_Pw_Pc_P\\
 &\ge\frac13+\frac{\sqrt{1-\tau}}{2\sqrt3}.
\end{align*}
This proves the profile for the entire critical locus.

Finally, for $a,b\ge0$, $a+b>0$, the critical normal family
\begin{equation}\label{eq:fqe-covering-equality-family}
 v=\frac{(a+ib,-a+ib,0,0)}{\sqrt{2(a^2+b^2)}},\qquad
 \tau=\frac{(a^2-b^2)^2}{(a^2+b^2)^2},\qquad
 \Phi=\frac13+\frac{ab}{\sqrt3(a^2+b^2)}
\end{equation}
attains equality. FP-4 checks every one of the $32$ ray overlaps against
this maximum. The parameters cover $0\le\tau\le1$, proving pointwise
optimality and the conditional covering radius. The endpoint $a=b$
is a cluster state; $b=0$ is GHZ up to LU.
